% =========================================================================
% background.tex
% =========================================================================

\section{Background}
\label{sec:background}

This section reviews the four source papers whose methods are evaluated
at $\dhead = 64$.

% -----------------------------------------------------------------
\subsection{TurboQuant and PolarQuant: Rotation-Based Scalar Quantization}
\label{sec:bg:turboquant}

TurboQuant~\cite{zandieh2025turboquant} compresses KV cache vectors
through a two-stage pipeline.
In the first stage, each $d$-dimensional vector $\mathbf{x}$ is
multiplied by a random orthogonal matrix $\Pi$, producing
$\mathbf{y} = \Pi\mathbf{x}$.
By the properties of random rotations on the unit sphere, each squared
coordinate $y_i^2$ follows a
$\mathrm{Beta}(1/2,\, d/2 - 1/2)$ distribution.
For large $d$, the CLT implies that each $y_i$ is
approximately Gaussian with known variance, enabling the use of
pre-computed, data-oblivious Lloyd-Max scalar
quantizers~\cite{lloyd1982least}.
TurboQuant implements the rotation via either a QR decomposition of a
random Gaussian matrix or a WHT composed with a random sign flip.
The WHT variant admits $O(d \log d)$ computation versus $O(d^2)$ for
full QR, though this advantage requires fused kernel support to realize
in practice.

In the second stage (TurboQuantProd), a one-bit QJL
transform~\cite{zandieh2024qjl} is applied to the quantization residual.
This corrects the systematic bias that MSE-optimal scalar quantizers
introduce when estimating inner products, yielding an unbiased estimator
of $\langle \mathbf{x}, \mathbf{x}' \rangle$ from the quantized
representations.

PolarQuant~\cite{han2025polarquant} pursues a related strategy: after
random preconditioning, it converts coordinate pairs to polar form
$(r, \theta)$ and quantizes the resulting angles, whose distribution
is analytically computable and tightly concentrated.
Both TurboQuant and PolarQuant report experimental validation
exclusively on models with $\dhead \geq 128$.

% -----------------------------------------------------------------
\subsection{QJL: One-Bit Quantized Johnson-Lindenstrauss Transform}
\label{sec:bg:qjl}

The QJL transform~\cite{zandieh2024qjl} projects a vector
$\mathbf{x} \in \mathbb{R}^d$ using a random matrix
$\mathbf{S} \in \mathbb{R}^{m \times d}$ and retains only the sign
bits $\mathrm{sign}(\mathbf{S}\mathbf{x})$.
By the Johnson-Lindenstrauss lemma~\cite{johnson1984extensions},
distances are approximately preserved under random projection.
The key insight of QJL is an asymmetric inner-product estimator:
applying the quantized projection to one vector and a standard
(unquantized) JL projection to the query yields an unbiased estimate
of their inner product with zero storage overhead for quantization
constants.

When used as a residual correction atop TurboQuant's MSE quantizer,
QJL adds one bit per coordinate.
At $\dhead = 128$, the variance of this estimator is small relative to
the signal.
However, the variance scales as $O(1/d)$; at $\dhead = 64$, the
variance doubles relative to $\dhead = 128$, and the exponential
sensitivity of softmax attention amplifies these fluctuations.
This variance-versus-bias tradeoff is the mechanism behind the failure
mode documented in Section~\ref{sec:method:adaptive}.

% -----------------------------------------------------------------
\subsection{Memory Caching: Segment Checkpointing with Growing Memory}
\label{sec:bg:mc}

Memory Caching (MC)~\cite{behrouz2026mc} augments recurrent models by
periodically checkpointing their hidden states.
The input sequence is divided into segments of size $S$; at each segment
boundary, the current hidden state is appended to a growing cache.
During attention, the model attends over both the recurrent state and
all cached checkpoints, interpolating between the $O(L)$ complexity of
RNNs and the $O(L^2)$ complexity of transformers.

Behrouz et al.\ propose four MC variants:
\begin{itemize}[noitemsep,topsep=1pt,parsep=0pt,partopsep=0pt]
\item \textbf{Residual Memory (RM):} Caches the change in hidden state
  between segments, so that older memories decay naturally.
\item \textbf{Gated Residual Memory (GRM):} Applies a similarity-based
  gate that modulates the update magnitude, preventing catastrophic
  forgetting of relevant past states.
\item \textbf{Memory Soup:} Uses softmax-weighted aggregation across
  multiple memory channels.
\item \textbf{Sparse Selective Caching (SSC):} Caches only top-$k$
  entries ranked by query-key similarity.
\end{itemize}

The original paper trains models of 760M to 1.3B parameters from
scratch on 30--100 billion tokens and evaluates on Multi-Query
Associative Recall (MQAR) and language modeling.
This study applies MC at inference time to a pre-trained GPT-2 Medium
model and trains a 15M-parameter variant from scratch on WikiText-103.
